\documentclass[11pt]{article}
\usepackage[margin=1in]{geometry}
\usepackage{amsmath,amssymb,amsthm}
\usepackage{graphicx}
\usepackage[colorlinks=true,linkcolor=blue,citecolor=blue,urlcolor=blue]{hyperref}
\usepackage{microtype}

\newtheorem{theorem}{Theorem}
\newcommand{\N}{\mathbb N}

\title{A sublacunary interpolation set for distal sequences}
\author{Run \texttt{fa\_banach\_001}, lane 18
  (\texttt{agent\_lane\_18})}
\date{13 August 2026}

\begin{document}
\maketitle

\begin{abstract}
Question 6.4 of Le's paper asks whether no sublacunary set can be an
interpolation set for distal sequences.  A result already implicit in
Pavlov's 2008 construction and stated explicitly in a 2024 follow-up gives
a negative answer.  The square set $\{n^2:n\ge1\}$ is sublacunary and is an
interpolation set for distal systems.  We record the two one-line growth
computations, identify the precise literature implication, and distinguish
this solved question from the source's two questions that remain open.
\end{abstract}

\section{The exact question}

An increasing set $E=\{r_1<r_2<\cdots\}\subset\N$ is
\emph{sublacunary} in the source if
\[
 \lim_{n\to\infty}\frac{\log r_n}{n}=0.
\]
It is an \emph{interpolation set for distal sequences} if every bounded
$f:E\to\mathbb C$ extends to a sequence $F(T^nx)$ arising from a continuous
function on a compact distal dynamical system.  The source asks:
\begin{quote}
\emph{Question 6.4. Is it true that no sublacunary set is an interpolation
set for distal sequences?}
\end{quote}
The exact source page, including the adjacent finite-union question, is
reproduced below.
\begin{center}
 \includegraphics[width=0.86\textwidth]{figures/open_problem_crop.png}
\end{center}

\section{The literature criterion}

Pavlov's Corollary 5.1 \cite{Pavlov} constructs arbitrary prescribed circle
coordinates along sequences whose gaps grow sufficiently rapidly, using
iterated skew products on tori.  These systems are distal extensions of an
irrational rotation.  Koutsogiannis, Le, Moreira, Pavlov, and Richter state
the exact interpolation consequence as follows \cite[Section 7]{KLMPR}:

\begin{theorem}[Pavlov criterion]\label{criterion}
Let $S=\{s_1<s_2<\cdots\}\subset\N$.  If
\begin{equation}\label{growth}
 \limsup_{n\to\infty}
 \frac{\log s_{n+1}}{\log(s_{n+1}-s_n)}<\infty,
\end{equation}
then $S$ is an interpolation set for distal systems.
\end{theorem}

The original circle-valued realization yields arbitrary bounded complex
values by normalization and a product construction.  Indeed, every
$z=re^{i\theta}$ in the closed unit disk is the average of the two unit
complex numbers
\[
 e^{i(\theta+\arccos r)}
 \quad\text{and}\quad
 e^{i(\theta-\arccos r)}.
\]
Products of distal systems are distal, and averaging their two continuous
coordinate functions gives the required complex extension.

\section{The square-set counterexample}

\begin{theorem}
The set
\[
 S=\{n^2:n\in\N\}
\]
is sublacunary and is an interpolation set for distal sequences.
Consequently Question 6.4 has a negative answer.
\end{theorem}

\begin{proof}
Put $s_n=n^2$.  The criterion \eqref{growth} applies because
\[
 \frac{\log s_{n+1}}{\log(s_{n+1}-s_n)}
 =\frac{2\log(n+1)}{\log(2n+1)}\longrightarrow2.
\]
Thus $S$ is a distal interpolation set.  On the other hand,
\[
 \frac{\log s_n}{n}=\frac{2\log n}{n}\longrightarrow0,
\]
so $S$ is sublacunary in the exact sense of the source.
\end{proof}

The same computation gives a continuum of elementary counterexamples:
$s_n=\lfloor n^{1+\varepsilon}\rfloor$ for any $\varepsilon>0$ (after
discarding repetitions, if needed) satisfies \eqref{growth} and is
sublacunary.

\section{Status of the adjacent questions}

This packet answers only the question for which the literature already gives
a complete resolution.  The 2024 follow-up explicitly preserves the other
two questions:
\begin{itemize}
\item The source's Question 6.3---whether a distal interpolation set remains
one after adjoining finitely many integers---reappears as Question 7.1 in
\cite{KLMPR}.  Proposition 7.2 reduces a positive answer to a splitting
property for sets of pointwise distal recurrence, and that property remains
open.
\item The strict hierarchy between $k$-step and $(k+1)$-step interpolation
sets is known at $k=1$ but reappears for $k\ge2$ as Question 8.2 in
\cite{KLMPR}.
\end{itemize}

Thus the correct classification is \emph{literature-implied answer}, not a
new result: Pavlov's theorem predates the target, and the later paper by Le
and collaborators makes its interpolation consequence explicit.

\begin{thebibliography}{9}
\bibitem{Le}
A. N. Le,
\emph{Sublacunary sets and interpolation sets for nilsequences},
Discrete Contin. Dyn. Syst. \textbf{42} (2022), 1855--1871;
arXiv:2103.13551.

\bibitem{Pavlov}
R. Pavlov,
\emph{Some counterexamples in topological dynamics},
Ergodic Theory Dynam. Systems \textbf{28} (2008), 1291--1322,
especially Corollary 5.1.

\bibitem{KLMPR}
A. Koutsogiannis, A. N. Le, J. Moreira, R. Pavlov, and F. K. Richter,
\emph{Interpolation sets for dynamical systems},
Trans. Amer. Math. Soc. \textbf{378} (2025), 1373--1400;
arXiv:2401.15339.
\end{thebibliography}

\end{document}

